% Personal preparation companion. Main report reordered on 02/10/2026; numbers unchanged.
\documentclass[10pt,a4paper]{article}
\usepackage{fontspec}
\setmainfont{texgyretermes}[Extension=.otf,UprightFont=*-regular,BoldFont=*-bold,ItalicFont=*-italic,BoldItalicFont=*-bolditalic]
\usepackage[margin=1.9cm,headheight=14pt]{geometry}
\usepackage{amsmath,amssymb,booktabs,longtable,array,ragged2e,enumitem,microtype,xcolor,graphicx,fancyhdr,tikz}
\usetikzlibrary{arrows.meta,positioning}
\usepackage[hidelinks,unicode]{hyperref}
\definecolor{ink}{HTML}{183A45}\definecolor{teal}{HTML}{087F7A}\definecolor{pale}{HTML}{EAF4F2}\definecolor{gray}{HTML}{56656B}
\newcolumntype{P}[1]{>{\RaggedRight\arraybackslash}p{#1}}
\setlength{\parindent}{0pt}\setlength{\parskip}{6pt}\setlength{\emergencystretch}{2em}
\setlength{\tabcolsep}{4pt}\renewcommand{\arraystretch}{1.15}
\setlist{nosep,leftmargin=1.4em,topsep=4pt,itemsep=3pt}
\pagestyle{fancy}\fancyhf{}\fancyhead[L]{\small\color{gray}GHI CHÚ CHUẨN BỊ TRÌNH BÀY}\fancyhead[R]{\small\color{gray}03/10/2026}\fancyfoot[C]{\small\thepage}\renewcommand{\headrulewidth}{0.3pt}
\newcommand{\key}[1]{\par\smallskip\noindent\colorbox{pale}{\parbox{\dimexpr\linewidth-2\fboxsep}{#1}}\par\smallskip}
\newcommand{\say}[1]{\par\textbf{Có thể trình bày:}\ \textit{#1}\par}
\newcommand{\ask}[2]{\par\textbf{Nếu được hỏi: #1}\par #2\par}
\newcommand{\where}[1]{\par{\small\color{gray}Đọc cùng report chính: #1.}\par}
\newcommand{\transition}[1]{\par\textbf{Chuyển ý:}\ #1\par}
\newcommand{\example}{\textbf{Ví dụ minh họa, không phải kết quả benchmark.} }
\newcommand{\guidepage}[1]{\clearpage\section{#1}}
\newcommand{\layer}[2]{\clearpage\noindent{\color{teal}\small\bfseries #1}\par\vspace{2pt}{\Large\bfseries\color{ink}#2}\par\vspace{6pt}\hrule\vspace{8pt}}
\hypersetup{pdftitle={Chuẩn bị trình bày nghiên cứu: related works, kiến trúc và benchmark},pdfauthor={}}
% Geo-I tables are shared via geoi_configuration.tex and geoi_benchmark.tex.

\begin{document}
\begin{center}
{\LARGE\bfseries\color{ink}Chuẩn bị thuyết trình}\\[5pt]
{\large Related works → kiến trúc → benchmark}\\[3pt]
{\small Buổi 03/10/2026 · Đọc cùng report rút gọn · Geo-I là phương pháp chính}
\end{center}
\key{Mạch nói: papers/scenario/metrics → các lớp mô hình → kết quả đối chứng. Mở phụ lục khi cần giải thích A/B/C.}

% Generated by build_report.py; shared task-level coverage table.
\section{Related works: mức bao phủ S1–S10}
Khảo sát giữ cửa sổ \textbf{21/09/2023–21/09/2026}. Bảng đối chiếu từng paper với mười nhiệm vụ của ta; mức bao phủ phải đọc cùng giả định và giới hạn ở cột cuối.

S1: vị trí hiện tại; S2: nơi dừng; S3: đoạn đường đã đi; S4: liên kết danh tính giữa phiên; S5: cạnh đường kế tiếp; S6: đích chưa tới; S7: nội dung/ý định truy vấn; S8: định vị nhờ dữ liệu người đi cùng; S9: điểm đầu bị che; S10: điểm cuối bị che của chuyến đã hoàn tất.

\textbf{Fully (F):} paper trực tiếp xử lý mục tiêu suy luận của scenario, trong giả định của paper. \textbf{Partially (P):} chỉ xử lý một phần hoặc bối cảnh hẹp hơn. \textbf{CX:} chưa đủ nguồn để phân loại. \textbf{—:} chưa tìm thấy bằng chứng xử lý trực tiếp; không kết luận cơ chế chắc chắn thất bại. F/P là đối chiếu của ta, không phải nhãn do tác giả công bố, không có nghĩa bảo vệ 100\% hoặc đã vượt các ca A/B/C của ta.

\begingroup\small
\begin{longtable}{@{}P{3.387cm}P{1.742cm}P{2.613cm}P{8.614cm}@{}}
\toprule
\textbf{Paper} & \textbf{Fully} & \textbf{Partially / CX} & \textbf{Căn cứ và giới hạn}\\\midrule
\endfirsthead
\toprule
\textbf{Paper} & \textbf{Fully} & \textbf{Partially / CX} & \textbf{Căn cứ và giới hạn}\\\midrule
\endhead
\bottomrule\endfoot
TransProtect 2024 [\href{https://arxiv.org/html/2409.09495v1}{R1}] & S1, S3 & S2 & §3–5: suy vị trí từ chuỗi đã làm nhiễu, biết đường/traffic; chưa đánh giá riêng nơi dừng.\\
\addlinespace[3pt]
Semantic correlation 2026 [\href{https://link.springer.com/article/10.1007/s44443-026-00899-w}{R2}] & S1, S3 & S2 & §3–6: giữ dummy hợp lý theo thời gian/ngữ nghĩa; chưa chấm riêng nơi dừng hay ý định S7.\\
\addlinespace[3pt]
Fake queries 2026 [\href{https://link.springer.com/article/10.1007/s44443-025-00438-z}{R3}] & S1, S3 & S2 & §3–6: chèn truy vấn giả để khó nối tuyến; chưa đánh giá riêng tọa độ nơi dừng.\\
\addlinespace[3pt]
Road-aware PTPPM 2025 [\href{https://arxiv.org/html/2511.21020v1}{R4}] & S1, S3 & S2 & §III, VI: chống suy vị trí/quỹ đạo có tương quan; bước dự báo chưa là phép thử S5/S6.\\
\addlinespace[3pt]
CPCROK 2025 [\href{https://uhra.herts.ac.uk/id/eprint/25705/1/futureinternet-17-00165.pdf}{R5}] & — & S4 & §3–5: chống nối bí danh xe trong VANET; chưa tách danh tính người/thiết bị như S4.A/B/C.\\
\addlinespace[3pt]
DP-FETC 2025 [\href{https://www.aimspress.com/aimspress-data/era/2025/11/PDF/era-33-11-293.pdf}{R6}] & — & S3, S8, S9, S10 & SynFE/PreOD: sinh dữ liệu offline, đổi đầu/cuối của tuyến tương quan; không chấm định vị từ người đi cùng.\\
\addlinespace[3pt]
Improved PIR 2025 [\href{https://onlinelibrary.wiley.com/doi/abs/10.1002/cpe.70447}{R7}] & CX & CX: S1, S7 & Tóm tắt nêu bảo vệ vị trí/nội dung bằng mã hóa; thiếu toàn văn để đối chiếu quyền quan sát và suy ý định.\\
\addlinespace[3pt]
SoK 2024 [\href{https://petsymposium.org/popets/2024/popets-2024-0068.pdf}{R8}] & Không áp dụng & Không áp dụng & Khảo sát cách đánh giá/tấn công; không phải cơ chế bảo vệ để gán F/P.\\
\addlinespace[3pt]
PRISM 2025 [\href{https://www.sciencedirect.com/science/article/pii/S1084804525001687}{R9}] & CX & CX: S1, S3 & Tóm tắt: bảo vệ chia sẻ quỹ đạo; chưa đủ mô hình attacker. Dùng LSTM không xác nhận xử lý S5/S6.\\
\addlinespace[3pt]
Options to Action 2026 [\href{https://elathan.github.io/papers/chi26.pdf}{R10}] & Không áp dụng & Không áp dụng & Khảo sát sử dụng tính năng, có che đầu/cuối (S9/S10); không đo khả năng chống suy luận.\\
\addlinespace[3pt]
Forsch et al. 2023 [\href{https://link.springer.com/chapter/10.1007/978-3-031-35374-1_3}{R11}] & S9, S10 & — & §3.2: S-TT cắt cả hai đầu, xét nhiều tuyến cùng đích; bảo đảm theo tập tấn công/địa điểm cho trước, offline.\\
\addlinespace[3pt]
EV querying 2024 [\href{https://wrap.warwick.ac.uk/id/eprint/183198/1/WRAP-privacy-preserving-querying-mechanism-high-utility-electric-vehicles-2024.pdf}{R12}] & S1, S3 & S2 & Proposed Query Model: AGeoI + dummy chống định vị/nối tuyến online; chưa chấm riêng nơi dừng hoặc S4–S10.\\
\end{longtable}
\endgroup
Scenario không được liệt kê ở một hàng được hiểu là —; với R7/R9 là chưa xác minh, với R8/R10 là không áp dụng. Bảng cho thấy các cơ chế chuỗi chủ yếu xét S1/S3; S2 cần phép thử nơi dừng riêng, còn S9/S10 cần kiểm tra suy ngược từ phần tuyến vẫn công bố.


\clearpage\subsection{Từ related works đến bộ đo đề xuất}
\say{Em đã cập nhật related works trong cửa sổ ba năm. Bảng cho thấy paper nào trực tiếp xét nhiệm vụ của mình, paper nào chỉ xử lý một phần. Các cơ chế chuỗi chủ yếu xét S1/S3; nơi dừng và suy điểm đầu/cuối vẫn cần phép thử riêng. Fully ở đây là bao phủ nhiệm vụ trong giả định paper, chưa có nghĩa vượt mọi ca A/B/C của dataset.}

\begin{longtable}{@{}P{4.1cm}P{5.4cm}P{6.1cm}@{}}
\toprule Paper tiêu biểu & Metrics gốc & Điều rút ra\\\midrule
TransProtect 2024 & EIE: sai số suy luận; chênh chi phí hành trình & Gần mục tiêu riêng tư; chi phí hành trình chưa cho biết có tìm đúng POI.\\
Semantic / Fake queries 2026 & ASR, DER hoặc số đường khó phân biệt; thời gian sinh & ASR là mức đạt ẩn danh theo paper, không tự đổi thành tỷ lệ đối thủ đoán đúng.\\
DP-FETC 2025 & MI: lượng thông tin; JSD: độ khác phân bố; bảo đảm DP & Giữ phân bố dữ liệu chưa có nghĩa giữ chất lượng dịch vụ hay che tốt điểm đầu/cuối.\\
CPCROK / PRISM 2025 & Khớp quỹ đạo, số giả; LPI, TDD, DC & Tác vụ khác hoặc thiếu định nghĩa đã xác minh; không so điểm chỉ vì tên chỉ số gần nhau.\\
SoK 2024 & Kiểm tra tấn công và chất lượng theo tác vụ & Cơ sở để chấm trên cùng nhiệm vụ; không phải model đối chứng.\\
\bottomrule
\end{longtable}
Đầu ra khác nhau (một điểm thay thế, một quỹ đạo giả, tập dummy hay query chèn) khiến metric dựa vào từng dummy không dùng chung được. Khác threat model, nhiệm vụ và dataset cũng ảnh hưởng. Ta chấm toàn bộ dữ liệu được phép thấy để suy cùng đáp án. Bảng metrics đủ 12 nguồn ở report. DLS/RDG là mốc nền cũ; AnotherMe online 11/09/2023, ngoài cửa sổ dù số tạp chí ghi 2024.

\textbf{Ba câu hỏi dẫn tới bộ đo đề xuất:}
\begin{itemize}
\item \textbf{Đối thủ đoán đúng tới đâu?} Hit50/100/200 là tỷ lệ đoán cách đáp án không quá 50/100/200 m; thấp hơn tốt hơn. MAE là sai số trung bình; median là trung vị; p90 là ngưỡng chứa 90\% sai số mẫu. Đọc cả Hit và sai số để thấy những lần vẫn lộ rất gần.
\item \textbf{Người dùng nhận được gì?} Recall@5 đo phần địa điểm chuẩn tìm lại được. Ví dụ, đúng 4/5 là 80\%. So với 5 POI chuẩn theo GPS thật, cùng trạng thái máy chủ; mỗi ca cần ≥90\%. “Trả đủ” đo số lượng, $\Delta D$ đo khoảng cách đường tăng thêm.
\item \textbf{Chi phí là bao nhiêu?} Đếm byte gửi + nhận trên mỗi lần dùng dịch vụ thật, gồm truy vấn chèn và tải ngoài chuyến. Độ trễ p50/p95 là trung vị/mức 95\%; hiện chưa đo toàn dịch vụ, chưa có điểm tổng Q thực nghiệm.
\end{itemize}
\say{Đóng góp không phải tạo ra MAE hay Recall mới. Em chọn và áp dụng chung các chỉ số đã có: đối thủ đoán cùng đáp án, người dùng dùng cùng dịch vụ và chi phí được tính đầy đủ.}
\transition{Từ các mục tiêu đó, em trình bày các lớp của mô hình bảo vệ.}

\clearpage\section{Kiến trúc mô hình: Geo-I và các lớp cơ chế}
\begin{center}\includegraphics[width=\linewidth]{figures/architecture_model.pdf}\end{center}
\textbf{Đọc hình:} input ở bên trái; ba layer bên trong khung xanh lần lượt bảo vệ Geo-I, xây dựng ngữ cảnh và sinh dummy. Output công khai là tập dummy và thời điểm phát; output riêng là top-5 POI tại thiết bị. GPS thật tạo neo và xếp hạng cuối tại thiết bị, không đi vào bộ chọn dummy. Khung xanh là mô hình của ta, gồm cả layer 4 boundary ở hàng dưới (S9 trước lõi, S10 sau lõi).
\say{Nền tảng của mô hình là Geo-I. Từ GPS thật, em lấy ngẫu nhiên một vị trí đại diện đã được bảo vệ, gọi là neo. Các bước sau dùng lịch sử neo, mạng đường và POI công khai để sinh tập điểm giả. Em bổ sung xử lý biên cho điểm đầu/cuối, thay vì coi nhiễu từng vị trí là đủ.}
\begin{enumerate}
\item \textbf{Geo-I / REM — S1:} lấy neo bằng cơ chế mũ REM: vị trí càng xa GPS thì xác suất chọn càng thấp. Tập ứng viên đường là công khai, khoảng cách dùng Euclid. Đây là nền bảo vệ từng vị trí.
\item \textbf{Tái dùng neo, ngân sách và pacing — S2/S3:} kiểm tra khoảng cách có nhiễu để giữ neo cũ hoặc tạo neo mới; cả hai bước tiêu ngân sách riêng tư. Tái dùng giúp giảm nhiều mẫu độc lập khi dừng. Ngân sách giới hạn rò rỉ tích lũy theo phiên (composition). Hết ngân sách thì chỉ xử lý neo cũ. Pacing giãn đọc GPS, chưa che giờ chuyến.
\item \textbf{Belief + road-aware + POI-aware — S2/S3 và chất lượng dịch vụ:} belief là phân bố ước lượng xe có thể ở đâu từ neo đã bảo vệ. Road-aware giữ dummy đi được theo hướng làn, luật rẽ và thời gian; POI-aware chọn K điểm phủ địa điểm bổ sung nhau. Exchange đổi ứng viên; bản GeoI-Slack cho phép độ lệch nhỏ trong objective khi chọn. Switching dừng/đi là biến thể trước, không phải cấu hình đang báo số. Bộ chọn không đọc tốc độ thật hay đích tương lai.
\item \textbf{BR-Boundary — S9/S10:} bỏ GPS của h giây đầu \emph{trước lõi}, tránh lưu dấu điểm đầu trong neo. Giữ bản tin đã bảo vệ trong bộ đệm Δ giây \emph{sau lõi} rồi mới phát; chuyến kết thúc thì hủy phần còn lại. Cách này không cần biết trước đích nhưng tăng độ trễ.
\end{enumerate}
\key{Lõi S1–S3: Geo-I + xử lý thời gian + mạng đường/POI. Mở rộng S9/S10: bỏ đầu trước lõi và buffer đuôi sau lõi. Các lớp boundary đã có kiểm tra code/tích hợp, chưa có benchmark toàn bộ mô hình với dịch vụ.}
\ask{Road-aware có tự tăng bảo đảm Geo-I không?}{Không. Nó làm dummy khả thi trên đường và hỗ trợ dịch vụ. Bảo đảm riêng tư đến từ cơ chế neo và cách cộng ngân sách. Các bước chỉ dùng dữ liệu đã bảo vệ giữ bảo đảm đó theo tính chất hậu xử lý. Sức chống đối thủ biết đường vẫn phải đo. K điểm giả không tự tạo k-anonymity.}
\transition{Tiếp theo, em nêu tham số của hai cấu hình Geo-I và dùng số liệu để đánh giá từng đóng góp.}

\clearpage
% Generated by build_report.py; shared Geo-I content.
\subsection{Định lượng hai cấu hình Geo-I và đóng góp của mô hình}
Hai tên dễ nhớ dưới đây đều thuộc \textbf{BR-Dummy trên nền Geo-I}, dùng cùng ngân sách riêng tư. GeoI-Paced là mốc so sánh nội bộ; GeoI-Slack là cấu hình làm việc hiện tại. Chúng khác cách chọn dummy, không phải hai mức $\varepsilon$ ưu tiên privacy/utility khác nhau.

\begingroup\small
\begin{longtable}{@{}P{3.938cm}P{6.350cm}P{6.350cm}@{}}
\toprule
\textbf{Tham số / component} & \textbf{GeoI-Paced} & \textbf{GeoI-Slack}\\\midrule
\endfirsthead
\toprule
\textbf{Tham số / component} & \textbf{GeoI-Paced} & \textbf{GeoI-Slack}\\\midrule
\endhead
\bottomrule\endfoot
Mã thực nghiệm & response\_\allowbreak{}paced / epoch\_\allowbreak{}cache & response\_\allowbreak{}paced\_\allowbreak{}slack03 / epoch\_\allowbreak{}cache\\
\addlinespace[3pt]
Neo riêng tư & REM; $\varepsilon$\_\allowbreak{}test = $\varepsilon$\_\allowbreak{}release = 0,01 /m; θ = 200 m & Cùng tham số neo\\
\addlinespace[3pt]
Ngân sách phiên & 23 đơn vị $\times$ 0,01; cận phiên 0,23 /m & Cùng cận phiên 0,23 /m\\
\addlinespace[3pt]
Pacing đọc GPS & Cách nhau ít nhất 60 giây & Cùng nhịp đọc\\
\addlinespace[3pt]
Road / belief / POI & Belief từ neo; làn có hướng, luật rẽ; chọn tập phủ POI & Cùng nền; thêm bước di chuyển có slack\\
\addlinespace[3pt]
Slack trong bộ chọn & Không nới objective & Cho phép mất tối đa 0,03 objective để tiến về mục tiêu; không phải giảm Recall 3\%\\
\addlinespace[3pt]
Đầu ra / dịch vụ & K=5 tọa độ; 6 loại POI; server trả top-10 mỗi loại/điểm & Cùng K=5, L=10\\
\addlinespace[3pt]
Cache và kết quả riêng & Hợp phản hồi cùng epoch 60 giây; GPS thật xếp hạng top-5 tại thiết bị & Cùng cache; không thay query\\
\end{longtable}
\endgroup
Lần đầu tạo neo chi 1 đơn vị; trước bước sau dự trữ 2 đơn vị. Test tái dùng chi 1, test + tạo neo mới chi 2. Giữa hai lần đọc hoặc khi hết ngân sách, mô hình chỉ dùng trạng thái đã bảo vệ. Nhịp đọc GPS, nhịp query dịch vụ và epoch hiệu lực phản hồi là ba khái niệm riêng; pacing không che giờ bắt đầu/kết thúc chuyến.

Trong phân tích lý tưởng với đồng hồ cố định, cận tỷ số xác suất của transcript là exp(0,23 $\times$ D$\infty$), với D$\infty$ là độ lệch GPS lớn nhất giữa hai chuỗi. Đây là cận theo khoảng cách, không phải xác suất “an toàn 23\%”. Hậu xử lý từ neo giữ cận này; mã số thực chưa có chứng nhận pure-DP chính xác.

\subsection{Đóng góp cần trình bày như thế nào?}
\textbf{Đề xuất:} từ lịch sử Geo-I đã bảo vệ, sinh tập truy vấn đi được trên mạng làn và phủ POI hữu ích trong ngân sách hữu hạn. Giá trị nằm ở sự phối hợp cơ chế theo nhiệm vụ dịch vụ: giảm suy luận gần, giữ khả năng tìm đúng địa điểm và kiểm soát tải truy vấn. Geo-I, tái dùng neo và dummy đã có tiền lệ; paper EV [\href{https://wrap.warwick.ac.uk/id/eprint/183198/1/WRAP-privacy-preserving-querying-mechanism-high-utility-electric-vehicles-2024.pdf}{R12}] đã kết hợp AGeoI + dummy, nên riêng tổ hợp hai từ này chưa là điểm mới.

\textbf{So với đối chứng:} DLS chú trọng độ khó phân biệt giữa điểm thật/dummy; các hướng TransProtect/Semantic chú trọng tính hợp lý của vị trí/chuỗi. Ta bổ sung việc chọn cả tập truy vấn theo hợp POI từ belief đã bảo vệ, thay vì dùng độ giống quỹ đạo làm đại diện cho utility. Road-aware tạo tính khả thi; privacy vẫn phải kiểm tra bằng attacker. Với RDG và Fake-query, đây mới là khác biệt thiết kế: chưa có benchmark Geo-I cùng protocol để chứng minh ưu thế số điểm.

\key{Lập luận đóng góp phải có cả cơ chế và bằng chứng: so trực tiếp bản Geo-I v2 với adapter để thấy đánh đổi; dùng phép bỏ/thêm component trên bản mới để kiểm tra lợi ích dịch vụ. Không gộp hai phiên bản thành một kết quả.}


\say{Geo-I là nền tảng của phương pháp. Em không nhận Geo-I hay dummy là phát minh mới. Đóng góp em kiểm tra là cách phối hợp lịch sử đã bảo vệ, ngân sách, mạng làn và mục tiêu phủ POI để giữ chất lượng dịch vụ. Hai cấu hình cùng ngân sách; bản Slack thêm cơ chế di chuyển trong bộ chọn.}

\clearpage
% Generated by build_report.py; shared Geo-I content.
\section{Benchmark Geo-I: so trực tiếp với đối chứng}
\textbf{Phép thử v2 đã được kiểm chứng:} K=5, ba seed 81/82/83, 12 chuyến test riêng biệt; mỗi chuyến có năm cửa sổ S1/S2/S3/S9/S10, nên 60 cửa sổ không phải 60 mẫu độc lập. Bản BR tái dùng dùng nền Geo-I với cận phiên 0,24 /m trên đồ thị nút giao. Đây là phiên bản trước mạng làn/pacing/slack/cache hiện tại, không có phân ca A/B/C.

* DLS là bản thích nghi lên đồ thị; TransProtect và Semantic là adapter rút gọn trong benchmark v2, không tái lập Transformer/LSTM của paper. Không thay chúng bằng adapter mới rồi giữ số cũ. RDG/Fake-query chưa có kết quả so trực tiếp với Geo-I trong phép thử này. AnotherMe chỉ áp dụng S3: 11/12 thành công, không cùng mẫu số nên báo riêng, không đưa vào bảng chung.

\begingroup\small
\begin{longtable}{@{}P{3.364cm}P{2.486cm}P{2.486cm}P{2.486cm}P{2.486cm}P{2.486cm}@{}}
\toprule
\textbf{Phương pháp} & \textbf{S1} & \textbf{S2} & \textbf{S3} & \textbf{S9} & \textbf{S10}\\\midrule
\endfirsthead
\toprule
\textbf{Phương pháp} & \textbf{S1} & \textbf{S2} & \textbf{S3} & \textbf{S9} & \textbf{S10}\\\midrule
\endhead
\bottomrule\endfoot
GPS thật & 100,00\% / 13 & 100,00\% / 39 & 96,43\% / 23 & 41,67\% / 120 & 8,33\% / 630\\
\addlinespace[3pt]
DLS* & 58,33\% / 1.318 & 100,00\% / 39 & 31,21\% / 1.688 & 8,33\% / 1.582 & 0,00\% / 2.029\\
\addlinespace[3pt]
TransProtect* & 83,33\% / 84 & 66,67\% / 88 & 69,28\% / 99 & 25,00\% / 184 & 8,33\% / 680\\
\addlinespace[3pt]
Semantic* & 100,00\% / 11 & 100,00\% / 42 & 40,58\% / 257 & 33,33\% / 132 & 16,67\% / 731\\
\addlinespace[3pt]
Geo-I + dummy & 58,33\% / 94 & 36,11\% / 105 & 52,65\% / 144 & 8,33\% / 307 & 0,00\% / 636\\
\addlinespace[3pt]
Geo-I + dummy đường & 16,67\% / 237 & 53,70\% / 121 & 3,57\% / 1.909 & 41,67\% / 217 & 0,00\% / 2.682\\
\addlinespace[3pt]
\textbf{BR tái dùng (Geo-I, v2)} & 33,33\% / 299 & 22,22\% / 247 & 7,44\% / 776 & 0,00\% / 316 & 0,00\% / 2.315\\
\end{longtable}
\endgroup
Mỗi ô: \textbf{Hit100 (\%) $\downarrow$ / MAE (m) $\uparrow$}. MAE và Hit chọn attacker riêng từ tập chọn đối thủ; không phải kết quả của một đối thủ duy nhất tối ưu cả hai.

\begingroup\small
\begin{longtable}{@{}P{3.364cm}P{2.486cm}P{2.486cm}P{2.486cm}P{2.486cm}P{2.486cm}@{}}
\toprule
\textbf{Recall@5 $\uparrow$} & \textbf{S1} & \textbf{S2} & \textbf{S3} & \textbf{S9} & \textbf{S10}\\\midrule
\endfirsthead
\toprule
\textbf{Recall@5 $\uparrow$} & \textbf{S1} & \textbf{S2} & \textbf{S3} & \textbf{S9} & \textbf{S10}\\\midrule
\endhead
\bottomrule\endfoot
GPS thật & 100,00\% & 100,00\% & 100,00\% & 100,00\% & 100,00\%\\
\addlinespace[3pt]
DLS* & 100,00\% & 100,00\% & 100,00\% & 100,00\% & 100,00\%\\
\addlinespace[3pt]
TransProtect* & 97,78\% & 97,19\% & 97,93\% & 97,27\% & 98,52\%\\
\addlinespace[3pt]
Semantic* & 100,00\% & 100,00\% & 100,00\% & 100,00\% & 100,00\%\\
\addlinespace[3pt]
Geo-I + dummy & 99,17\% & 98,52\% & 97,73\% & 96,85\% & 97,06\%\\
\addlinespace[3pt]
Geo-I + dummy đường & 98,06\% & 97,44\% & 81,96\% & 96,99\% & 79,75\%\\
\addlinespace[3pt]
\textbf{BR tái dùng (Geo-I, v2)} & 95,83\% & 97,01\% & 87,67\% & 96,60\% & 88,33\%\\
\end{longtable}
\endgroup
\textbf{Bằng chứng hỗ trợ:} ở S2, BR tái dùng có Hit100 \textbf{22,22\%}, thấp hơn DLS/Semantic (100\%) và TransProtect (66,67\%), với Recall 97,01\%. Ở S3, Hit100 \textbf{7,44\%}, thấp hơn ba adapter (31,21–69,28\%) và Geo-I + dummy (52,65\%). Kết quả hỗ trợ thiết kế phối hợp trạng thái đã bảo vệ theo thời gian; chưa cô lập riêng tác động tái dùng vì tham số và cách sinh dummy cũng khác.

\textbf{Đánh đổi phải nói rõ:} BR không thắng mọi metric. Geo-I + dummy đường có Hit100 S3 thấp hơn (3,57\%) nhưng Recall thấp hơn (81,96\% so với 87,67\%); cả hai chưa đạt 90\%. DLS có MAE lớn hơn ở S1/S3/S9 dù Hit100 có thể cao hơn. BR cũng chưa đạt Recall 90\% ở S10 (88,33\%); S10 Hit100=0 của nhiều phương pháp nên không đủ phân biệt. S9/S10 chỉ suy tọa độ đầu/cuối SUMO, chưa chứng minh che nhà hay nơi làm việc.

\key{Kết luận v2: BR trên nền Geo-I giảm tỷ lệ suy đúng gần ở một số nhiệm vụ quan trọng, có đánh đổi utility. Đây là bằng chứng đóng góp của bản thích nghi trong benchmark, chưa chứng minh vượt các paper nguyên bản hoặc ưu thế ở cùng ngân sách.}


\clearpage
\subsection{Bản Geo-I hiện tại: 14 ca và kiểm tra từng component}
\textbf{GeoI-Slack:} 12 nhóm phát triển, 165 bản ghi từ 94 chuyến nguồn; S10 chỉ A/B. K=5, L=10, cận phiên 0,23 /m. Recall dưới đây dùng POI khả dụng mô phỏng ở p=0,8, ba thế giới khả dụng và hai lần chạy cơ chế. Privacy kế thừa từ đúng tọa độ công bố đã chấm bằng tập attacker hữu hạn: cache và trạng thái server độc lập với GPS không đổi query. Đây là dữ liệu phát triển, không phải xác nhận trên nhóm mới.

\begingroup\small
\begin{longtable}{@{}P{1.589cm}P{2.980cm}P{3.477cm}P{2.980cm}P{3.179cm}P{1.589cm}@{}}
\toprule
\textbf{Ca} & \textbf{GPS thật: Hit100} & \textbf{GeoI-Slack: Hit100 $\downarrow$} & \textbf{MAE (m) $\uparrow$} & \textbf{Recall@5 $\uparrow$} & \textbf{Nhóm}\\\midrule
\endfirsthead
\toprule
\textbf{Ca} & \textbf{GPS thật: Hit100} & \textbf{GeoI-Slack: Hit100 $\downarrow$} & \textbf{MAE (m) $\uparrow$} & \textbf{Recall@5 $\uparrow$} & \textbf{Nhóm}\\\midrule
\endhead
\bottomrule\endfoot
S1.A & 100,00\% & 8,33\% & 852 & 97,64\% & 12\\
\addlinespace[3pt]
S1.B & 100,00\% & 8,33\% & 832 & 95,90\% & 12\\
\addlinespace[3pt]
S1.C & 100,00\% & 0,00\% & 1.674 & 80,18\% & 12\\
\addlinespace[3pt]
S2.A & 100,00\% & 0,00\% & 910 & 99,86\% & 12\\
\addlinespace[3pt]
S2.B & 100,00\% & 4,17\% & 715 & 99,49\% & 12\\
\addlinespace[3pt]
S2.C & 100,00\% & 0,00\% & 997 & 94,59\% & 12\\
\addlinespace[3pt]
S3.A & 100,00\% & 3,09\% & 906 & 95,63\% & 12\\
\addlinespace[3pt]
S3.B & 100,00\% & 3,37\% & 873 & 95,97\% & 11\\
\addlinespace[3pt]
S3.C & 100,00\% & 2,19\% & 934 & 95,26\% & 12\\
\addlinespace[3pt]
S9.A & 16,67\% & 0,00\% & 607 & 95,69\% & 12\\
\addlinespace[3pt]
S9.B & 75,00\% & 2,50\% & 701 & 95,30\% & 10\\
\addlinespace[3pt]
S9.C & 0,00\% & 0,00\% & 514 & 94,95\% & 12\\
\addlinespace[3pt]
S10.A & 33,33\% & 0,00\% & 1.345 & 97,63\% & 12\\
\addlinespace[3pt]
S10.B & 16,67\% & 0,00\% & 1.200 & 96,42\% & 12\\
\end{longtable}
\endgroup
\begingroup\small
\begin{longtable}{@{}P{4.839cm}P{3.581cm}P{3.194cm}P{4.742cm}@{}}
\toprule
\textbf{Cấu hình} & \textbf{Recall@5 $\uparrow$} & \textbf{Ca $\ge$90\%} & \textbf{Byte / sự kiện $\downarrow$}\\\midrule
\endfirsthead
\toprule
\textbf{Cấu hình} & \textbf{Recall@5 $\uparrow$} & \textbf{Ca $\ge$90\%} & \textbf{Byte / sự kiện $\downarrow$}\\\midrule
\endhead
\bottomrule\endfoot
GeoI-Paced & 94,38\% & 13/14 & 993,7\\
\addlinespace[3pt]
\textbf{GeoI-Slack} & 95,44\% & 13/14 & 994,1\\
\addlinespace[3pt]
GeoI-Slack, bỏ cache & 95,19\% & 13/14 & 994,1\\
\addlinespace[3pt]
GPS thật & 100,00\% & 14/14 & 214,5\\
\end{longtable}
\endgroup
\textbf{Đóng góp có thể định lượng:} cache tăng Recall 0,24 điểm \%, CI 95\% [0,14; 0,34], không thêm query/byte trong schema này. Thêm slack tăng trung bình 1,06 điểm \%, CI [-0,13; 2,23] chứa 0: hướng cải thiện có triển vọng, chưa đủ kết luận chắc chắn. CI lấy mẫu lại ghép cặp 12 nhóm 10.000 lần, chưa hiệu chỉnh nhiều so sánh.

\textbf{Giới hạn còn thấy:} S1.C chỉ đạt 80,18\% Recall, nên cấu hình đạt 13/14 ca $\ge$90\%, không phải đạt mọi ca. S9.C có Hit100=0 cả với GPS thật, nên riêng chỉ số này chưa chứng minh cải thiện; Geo-I vẫn có Hit500=58,33\%: điểm đầu vẫn có thể bị khoanh vùng. Những số S9/S10 là của lõi Geo-I qua cửa sổ quan sát; không gán cho BR-Boundary, vì wrapper cắt đầu/buffer đuôi chưa được benchmark cùng dịch vụ này.

Recall gộp đều ca trong scenario rồi đều năm scenario. Byte gộp theo sự kiện, tính toàn bộ traffic của 94 chuyến được giữ, gồm yêu cầu + phản hồi JSON; chưa gồm HTTP/TLS hoặc latency. S10.B ở đây giữ nguồn cũ S10.C. Các số tổng/CI được tính lại đúng 14 ca, không chạy lại mô hình và không trộn với kết quả v2.



\clearpage
% Generated by build_report.py; shared by both documents.
\section{Phụ lục: đọc chi tiết 14 sample A/B/C}\label{sec:4}
A/B/C là các điều kiện cùng nhiệm vụ, không phải thứ tự độ khó. Mỗi ca chọn một bản ghi thật từ bộ minh họa 12 nhóm/264 chuyến; số bản ghi không phải số chuyến độc lập. \textbf{S10 chỉ có A/B} trong tài liệu; B ánh xạ mã nguồn cũ S10.C để giữ truy nguyên, không có ca C độc lập trong phạm vi hiện tại.

Chấm màu = mẫu trước bảo vệ; nét đứt = đường thật để chấm; sao đỏ = đáp án; trục thời gian tách các mẫu chồng nhau. Đối thủ chỉ nhận dữ liệu đã bảo vệ và thông tin phụ trợ được phép, không nhận các tọa độ/nhãn thật này. Bản đồ minh họa dùng nền SUMO với cùng mã kiểm tra tệp OSM; thiếu mạng gốc để xác minh trùng toàn bộ hình học. © OpenStreetMap contributors.

\subsection{S1 — suy ra vị trí tại một lần gửi}\label{sub:4.1}
Đối thủ cần suy ra tọa độ tại một lần gửi. A/B thay ràng buộc mạng đường; C thêm ngữ cảnh POI hiếm.

\noindent\begin{minipage}[t]{0.48\linewidth}\vspace{0pt}\includegraphics[width=\linewidth]{figures/case_S1_A.pdf}\end{minipage}\hfill\begin{minipage}[t]{0.49\linewidth}\vspace{0pt}\RaggedRight \textbf{S1.A / 12 bản ghi.} Một bản tin tại cạnh có $\ge$2 hướng đi tiếp hợp lệ: kiểm tra khi mạng đường còn nhiều khả năng.\par\smallskip \textbf{Bản ghi minh họa v2-r301-000.} u301\_\allowbreak{}00: 1 mẫu, chỉ số GPS FCD[333…333]\par\smallskip\textbf{Đọc sample.} Một lần lấy mẫu ở cạnh có 4 hướng đi tiếp. Sao đỏ trùng vị trí đầu vào; đối thủ phải chọn vị trí từ đầu ra đã bảo vệ, không được thấy sẵn sao này.\end{minipage}\par\vspace{3pt}
\noindent\begin{minipage}[t]{0.48\linewidth}\vspace{0pt}\includegraphics[width=\linewidth]{figures/case_S1_B.pdf}\end{minipage}\hfill\begin{minipage}[t]{0.49\linewidth}\vspace{0pt}\RaggedRight \textbf{S1.B / 12 bản ghi.} Một bản tin tại cạnh chỉ có 1 hướng đi tiếp: bản đồ tự loại bớt khả năng.\par\smallskip \textbf{Bản ghi minh họa v2-r301-001.} u301\_\allowbreak{}00: 1 mẫu, chỉ số GPS FCD[241…241]\par\smallskip\textbf{Đọc sample.} Vẫn chỉ một lần lấy mẫu, nhưng cạnh chỉ có một hướng đi tiếp. Bản đồ có thể loại bớt khả năng so với A; cần đo lợi thế này bằng đối thủ thực tế.\end{minipage}\par\vspace{3pt}
\noindent\begin{minipage}[t]{0.48\linewidth}\vspace{0pt}\includegraphics[width=\linewidth]{figures/case_S1_C.pdf}\end{minipage}\hfill\begin{minipage}[t]{0.49\linewidth}\vspace{0pt}\RaggedRight \textbf{S1.C / 12 bản ghi.} Cách POI $\le$150 m; loại POI chiếm $\le$5\% danh mục công khai. Hiếm theo loại địa điểm, không phải theo tần suất ghé thăm.\par\smallskip \textbf{Bản ghi minh họa v2-r301-029.} u301\_\allowbreak{}13: 1 mẫu, chỉ số GPS FCD[695…695]\par\smallskip\textbf{Đọc sample.} Điểm lấy mẫu cách POI pharmacy 12,03 m; loại này chiếm 3,11\% danh mục. Ngữ cảnh địa điểm hiếm là dấu hiệu bổ sung, không phải chuỗi di chuyển.\end{minipage}\par\vspace{3pt}

\clearpage
\subsection{S2 — suy ra nơi dừng qua nhiều lần gửi}\label{sub:4.2}
Đối thủ cần suy ra tọa độ nơi dừng. A/B thay thời lượng và nhịp quan sát; C kiểm tra rời đi rồi quay lại.

\noindent\begin{minipage}[t]{0.48\linewidth}\vspace{0pt}\includegraphics[width=\linewidth]{figures/case_S2_A.pdf}\end{minipage}\hfill\begin{minipage}[t]{0.49\linewidth}\vspace{0pt}\RaggedRight \textbf{S2.A / 12 bản ghi.} Dừng có chủ đích $\ge$20 giây; lấy truy vấn mỗi 5 giây.\par\smallskip \textbf{Bản ghi minh họa v2-r301-002.} u301\_\allowbreak{}05: 6 mẫu, chỉ số GPS FCD[54…79]\par\smallskip\textbf{Đọc sample.} 6 mẫu cùng một tọa độ, cách nhau 5 giây; đoạn dừng thật dài 29 giây. Sáu vạch thời gian bị chồng thành một chấm trên map.\end{minipage}\par\vspace{3pt}
\noindent\begin{minipage}[t]{0.48\linewidth}\vspace{0pt}\includegraphics[width=\linewidth]{figures/case_S2_B.pdf}\end{minipage}\hfill\begin{minipage}[t]{0.49\linewidth}\vspace{0pt}\RaggedRight \textbf{S2.B / 12 bản ghi.} Dừng $\ge$120 giây; lấy truy vấn mỗi 20 giây.\par\smallskip \textbf{Bản ghi minh họa v2-r301-003.} u301\_\allowbreak{}06: 9 mẫu, chỉ số GPS FCD[59…219]\par\smallskip\textbf{Đọc sample.} 9 mẫu cùng một tọa độ, cách nhau 20 giây; đoạn dừng dài 179 giây. Đối thủ có cửa sổ dài hơn để kết hợp bằng chứng, không chỉ một bản tin.\end{minipage}\par\vspace{3pt}
\noindent\begin{minipage}[t]{0.48\linewidth}\vspace{0pt}\includegraphics[width=\linewidth]{figures/case_S2_C.pdf}\end{minipage}\hfill\begin{minipage}[t]{0.49\linewidth}\vspace{0pt}\RaggedRight \textbf{S2.C / 12 bản ghi.} Hai lần dừng trên cùng làn, mỗi lần $\ge$20 giây; ở giữa thực sự di chuyển.\par\smallskip \textbf{Bản ghi minh họa v2-r301-004.} u301\_\allowbreak{}07: 18 mẫu, chỉ số GPS FCD[214…1193]\par\smallskip\textbf{Đọc sample.} 18 mẫu chia hai đợt, mỗi đợt 9 mẫu. Hai lần dừng dài 44 giây tại cùng tọa độ; khoảng trống giữa hai cụm thời gian là lúc xe rời đi rồi quay lại.\end{minipage}\par\vspace{3pt}

\clearpage
\subsection{S3 — tái dựng đoạn đường đã đi}\label{sub:4.3}
Đối thủ cần tái dựng các vị trí của đoạn đã đi. A là chuỗi đều; B ít nhánh; C quan sát thưa.

\noindent\begin{minipage}[t]{0.48\linewidth}\vspace{0pt}\includegraphics[width=\linewidth]{figures/case_S3_A.pdf}\end{minipage}\hfill\begin{minipage}[t]{0.49\linewidth}\vspace{0pt}\RaggedRight \textbf{S3.A / 12 bản ghi.} Đoạn di chuyển $\ge$60 giây qua $\ge$3 cạnh; gửi mỗi 20 giây.\par\smallskip \textbf{Bản ghi minh họa v2-r301-005.} u301\_\allowbreak{}00: 12 mẫu, chỉ số GPS FCD[361…581]\par\smallskip\textbf{Đọc sample.} 12 mẫu của u301\_\allowbreak{}00 cách nhau 20 giây. Chuỗi điểm cung cấp nhiều ràng buộc để nối tuyến; nét đứt là đường thật chỉ dùng chấm.\end{minipage}\par\vspace{3pt}
\noindent\begin{minipage}[t]{0.48\linewidth}\vspace{0pt}\includegraphics[width=\linewidth]{figures/case_S3_B.pdf}\end{minipage}\hfill\begin{minipage}[t]{0.49\linewidth}\vspace{0pt}\RaggedRight \textbf{S3.B / 11 bản ghi.} Ít nhất một nửa số cạnh được lấy mẫu chỉ có một hướng đi tiếp: hành lang ít lựa chọn.\par\smallskip \textbf{Bản ghi minh họa v2-r301-007.} u301\_\allowbreak{}00: 9 mẫu, chỉ số GPS FCD[91…251]\par\smallskip\textbf{Đọc sample.} 9 mẫu; 55,6\% cạnh được lấy mẫu chỉ có một hướng đi tiếp. Đường ít nhánh có thể khiến chuỗi dễ suy hơn dù từng vị trí được làm mờ.\end{minipage}\par\vspace{3pt}
\noindent\begin{minipage}[t]{0.48\linewidth}\vspace{0pt}\includegraphics[width=\linewidth]{figures/case_S3_C.pdf}\end{minipage}\hfill\begin{minipage}[t]{0.49\linewidth}\vspace{0pt}\RaggedRight \textbf{S3.C / 12 bản ghi.} Lấy mẫu đoạn di chuyển thưa hơn, mỗi 60 giây.\par\smallskip \textbf{Bản ghi minh họa v2-r301-006.} u301\_\allowbreak{}00: 4 mẫu, chỉ số GPS FCD[361…541]\par\smallskip\textbf{Đọc sample.} 4 mẫu của cùng đoạn trong A, cách nhau 60 giây. Khoảng trống lớn hơn nhưng đối thủ vẫn có thể dùng mạng đường để nối các khả năng.\end{minipage}\par\vspace{3pt}

\clearpage
\subsection{S9 — suy ra điểm xuất phát bị che}\label{sub:4.4}
Đối thủ cần suy điểm đầu bị giấu. A suy ngược một chuyến; B phân biệt hai nguồn nhập tuyến; C kết hợp chuyến lặp.

\noindent\begin{minipage}[t]{0.48\linewidth}\vspace{0pt}\includegraphics[width=\linewidth]{figures/case_S9_A.pdf}\end{minipage}\hfill\begin{minipage}[t]{0.49\linewidth}\vspace{0pt}\RaggedRight \textbf{S9.A / 12 bản ghi.} Cạnh xuất phát chỉ một lối ra. Che 60 giây đầu, lấy phần còn lại mỗi 20 giây.\par\smallskip \textbf{Bản ghi minh họa v2-r301-025.} u301\_\allowbreak{}00: 27 mẫu, chỉ số GPS FCD[60…580]\par\smallskip\textbf{Đọc sample.} Sao là FCD[0]; cửa sổ bắt đầu ở FCD[60] sau khi che 60 giây. Chỉ một lối ra có thể giúp suy ngược nguồn xuất phát từ phần còn công bố.\end{minipage}\par\vspace{3pt}
\noindent\begin{minipage}[t]{0.48\linewidth}\vspace{0pt}\includegraphics[width=\linewidth]{figures/case_S9_B.pdf}\end{minipage}\hfill\begin{minipage}[t]{0.49\linewidth}\vspace{0pt}\RaggedRight \textbf{S9.B / 11 bản ghi.} Hai chuyến khác điểm đầu nhưng nhập một đoạn chung; chỉ cấp từ chỗ nhập tuyến.\par\smallskip \textbf{Bản ghi minh họa v2-r301-026.} u301\_\allowbreak{}00: 28 mẫu, chỉ số GPS FCD[34…574]; u301\_\allowbreak{}04: 28 mẫu, chỉ số GPS FCD[19…559]\par\smallskip\textbf{Đọc sample.} Hai điểm đầu cách 281,76 m rồi nhập tuyến. Chỉ cấp dữ liệu sau nhập; hai sao là hai đáp án bị giấu mà đối thủ phải phân biệt.\end{minipage}\par\vspace{3pt}
\noindent\begin{minipage}[t]{0.48\linewidth}\vspace{0pt}\includegraphics[width=\linewidth]{figures/case_S9_C.pdf}\end{minipage}\hfill\begin{minipage}[t]{0.49\linewidth}\vspace{0pt}\RaggedRight \textbf{S9.C / 12 bản ghi.} Nhiều chuyến có cùng điểm xuất phát dự kiến; che 60 giây đầu từng chuyến.\par\smallskip \textbf{Bản ghi minh họa v2-r301-027.} u301\_\allowbreak{}00: 27 mẫu, chỉ số GPS FCD[60…580]; u301\_\allowbreak{}09: 27 mẫu, chỉ số GPS FCD[60…580]\par\smallskip\textbf{Đọc sample.} u301\_\allowbreak{}00/u301\_\allowbreak{}09 lặp cùng điểm đầu dự kiến; mỗi phiên che 60 giây đầu. Nhiều cửa sổ có thể bổ sung bằng chứng về cùng nguồn xuất phát.\end{minipage}\par\vspace{3pt}

\clearpage
\subsection{S10 — suy ra điểm kết thúc bị che}\label{sub:4.5}
Đối thủ cần suy điểm cuối của chuyến đã hoàn tất. A dùng một chuyến; B kết hợp nhiều chuyến cùng đích. Đây khác dự đoán đích tương lai ở S6.

\noindent\begin{minipage}[t]{0.58\linewidth}\vspace{0pt}\includegraphics[width=\linewidth]{figures/case_S10_A.pdf}\end{minipage}\hfill\begin{minipage}[t]{0.39\linewidth}\vspace{0pt}\RaggedRight \textbf{S10.A / 11 bản ghi.} Cạnh kết thúc chỉ một lối vào. Che 60 giây cuối của chuyến đã hoàn tất.\par\smallskip \textbf{Bản ghi minh họa v2-r302-025.} u302\_\allowbreak{}08: 38 mẫu, chỉ số GPS FCD[0…740]\par\smallskip\textbf{Đọc sample.} Sao là FCD cuối [817] của u302\_\allowbreak{}08; che 60 giây cuối. Đường chỉ một lối vào có thể giúp suy đoạn kết thúc từ phần trước còn thấy.\end{minipage}\par\vspace{3pt}
\noindent\begin{minipage}[t]{0.58\linewidth}\vspace{0pt}\includegraphics[width=\linewidth]{figures/case_S10_B.pdf}\end{minipage}\hfill\begin{minipage}[t]{0.39\linewidth}\vspace{0pt}\RaggedRight \textbf{S10.B / 12 bản ghi.} Nhiều chuyến cùng điểm kết thúc dự kiến; che 60 giây cuối từng chuyến.\par\smallskip \textbf{Bản ghi minh họa v2-r301-028.} u301\_\allowbreak{}00: 27 mẫu, chỉ số GPS FCD[0…520]; u301\_\allowbreak{}09: 27 mẫu, chỉ số GPS FCD[0…520]\par\smallskip\textbf{Đọc sample.} u301\_\allowbreak{}00/u301\_\allowbreak{}09 lặp cùng đích dự kiến và cùng che 60 giây cuối. Đối thủ có thể kết hợp nhiều phiên để khoanh đích, chưa đủ gọi đó là nơi làm việc.\end{minipage}\par\vspace{3pt}
S9/S10 chấm tọa độ đầu/cuối; chưa gán nhãn nhà/nơi làm việc. Xem bản đồ phóng to của 29 ca của khung đầy đủ ở sample\_\allowbreak{}maps.html; tọa độ và bản ghi ở data\_\allowbreak{}samples.json, số đếm ở scenario\_\allowbreak{}inventory.csv.


\clearpage

\end{document}
